% =======================================================
% 1. INTRODUCTION
% =======================================================
\section{Introduction}
\label{sec:introduction}

The end of Dennard scaling and the transition to many-core and exascale systems \cite{sutter2005,flynn2007,dongarra2023} have made thread-level parallelism the dominant source of performance growth. A large body of scientific and industrial software, however, is still written as sequential loops. Converting it to OpenMP is an expert task: the programmer must establish, loop by loop, that no iteration reads or writes data that another iteration writes, and a single overlooked loop-carried dependence turns a performance optimization into a silent data race that compiles, often passes a small test, and is wrong.

\subsection{The automation gap and the loop-level focus}
Two families of tools try to automate this step. \emph{Compiler-based} parallelizers, from classical dependence tests \cite{banerjee1988,pugh1992} and polyhedral optimizers \cite{feautrier1991,bondhugula2008,grosser2012} to the auto-parallelization passes of production compilers \cite{gccparloops}, are sound by construction but conservative: whenever a dependence cannot be disproven at compile time---pointer parameters that alias analysis cannot separate, indirect subscripts such as \Code{A[B[i]]}, flattened multi-dimensional indices, conditionally updated scalars---the loop is left sequential. \emph{Learning-based} approaches, from transformer classifiers trained on OpenMP corpora \cite{harel2023pragformer,kadosh2023ompify,kadosh2024ompar} to general-purpose LLMs prompted with source code \cite{nichols2024pareval,mahmud2025autoparllm}, recognise algorithmic intent that compilers miss but provide no correctness guarantee; a model that confidently declares a loop with a cross-iteration dependence ``parallel'' produces exactly the kind of error that is hardest to find \cite{pearce2022asleep}.

We deliberately restrict the problem to \textbf{loop-level parallelization}: deciding, for one loop at a time, whether a single \Code{\#pragma omp parallel for} directive with data-sharing, reduction, linear and collapse clauses can be placed before it without changing the program's result. This is the decision a developer or a tool faces most often when modernizing serial code, it has a precise legality definition in terms of the OpenMP specification \cite{openmp52}, and it can be checked by execution. The restriction is what makes it possible to combine a sound but incomplete compiler analysis with a complete but unsound model \emph{without giving up soundness}: \SystemName{} lets the compiler decide every loop it can, vetoes every loop with a known dependence, and asks the model only about the rest.

\subsection{Relation to our preliminary version}
An earlier version of this work evaluated the same idea on 200 C++ files and reported an F1-score of 96.6\% for an LLVM-only configuration and 5.7\% for an LLM-only baseline. While re-implementing the system we found that this evaluation was not independent: its ground-truth labels were derived from the output of the same LLVM pass being evaluated, the pass counted conditional branches rather than natural loops, the prompt ablation was simulated, the 99\% cache hit rate was measured by re-analysing an unchanged corpus, and the confidence threshold was tuned on the test data. None of those numbers is reused here. All results in this article were produced from scratch with a re-implemented analysis, benchmarks whose labels are validated by execution, and real model calls, on one Apple M1 Pro machine; every number in the text and tables is generated by scripts from logged raw data (Section~\ref{sec:artifact}). The 200-file corpus of the preliminary study is retained as an application-scale corpus without ground truth (Section~\ref{sec:rq-field}).

\subsection{Contributions}
\begin{enumerate}
    \item \textbf{HF-Static}, an LLVM~21 pass that classifies every natural loop as \textsc{Parallel}, \textsc{Serial} or \textsc{Unknown} from LoopInfo, ScalarEvolution, DependenceAnalysis and the induction/recurrence descriptors, and reports \emph{why} a loop is undecided (Section~\ref{sec:static}).
    \item A \textbf{hybrid decision procedure} with a static veto and a call-site alias guard that queries an LLM only for \textsc{Unknown} loops, injects the compiler facts into the prompt, and gates the answer with a threshold $\tau$ calibrated on held-out kernels (Section~\ref{sec:method}).
    \item \textbf{Execution-validated loop benchmarks}: HF-LoopBench (\nBenchLoops{} loops, \nBenchDomains{} domains) and HF-LoopBench-X (\xLoops{} loops written after the study was frozen to stress loop structure), in which every parallel label is verified under 12 OpenMP configurations and every canonical serial label is checked for an observable race (Section~\ref{sec:bench}).
    \item An \textbf{empirical study} against GCC auto-parallelization, LLVM Polly and LLM-only classification with three open-weight models, covering threshold calibration, per-class and per-loop-structure analysis, statistical tests, a prompt ablation with real API calls, a replayed evaluation of answer caching, measured token cost under free-tier rate limits, run-time speedups, and an application-scale corpus (Sections~\ref{sec:setup}--\ref{sec:cases}).
\end{enumerate}
